\documentclass[12pt, a4paper]{article}

% ── Packages ──────────────────────────────────────────────────────────────────
\usepackage[margin=1in]{geometry}
\usepackage[T1]{fontenc}
\usepackage[utf8]{inputenc}
\usepackage{lmodern}
\usepackage{microtype}
\usepackage{xcolor}
\usepackage{titlesec}
\usepackage{booktabs}
\usepackage{tabularx}
\usepackage{array}
\usepackage{colortbl}
\usepackage{enumitem}
\usepackage{listings}
\usepackage{hyperref}
\usepackage{fancyhdr}
\usepackage{tcolorbox}
\usepackage{setspace}
\usepackage{amsmath}
\tcbuselibrary{skins, breakable}

% ── Colors ────────────────────────────────────────────────────────────────────
\definecolor{primaryblue}{HTML}{1A73E8}
\definecolor{darkblue}{HTML}{0D47A1}
\definecolor{lightgray}{HTML}{F5F5F5}
\definecolor{codebg}{HTML}{1E1E2E}
\definecolor{codetext}{HTML}{CDD6F4}
\definecolor{accent}{HTML}{E53935}
\definecolor{mysqlorange}{HTML}{E8710A}
\definecolor{tableheader}{HTML}{1A73E8}
\definecolor{successgreen}{HTML}{1B8A3E}
\definecolor{lightgreen}{HTML}{F0FFF4}

% ── Hyperref Setup ────────────────────────────────────────────────────────────
\hypersetup{
    colorlinks=true,
    linkcolor=primaryblue,
    urlcolor=primaryblue,
    pdftitle={GitDB: Implementation Plan},
    pdfauthor={Group Project - DBMS}
}

% ── Section Formatting ────────────────────────────────────────────────────────
\titleformat{\section}
  {\large\bfseries\color{darkblue}}
  {\thesection.}{0.5em}{}[\color{primaryblue}\titlerule]

\titleformat{\subsection}
  {\normalsize\bfseries\color{primaryblue}}
  {\thesubsection.}{0.5em}{}

\titleformat{\subsubsection}
  {\normalsize\bfseries\color{primaryblue!80}}
  {\thesubsubsection.}{0.5em}{}

\titlespacing*{\section}{0pt}{1.2em}{0.6em}
\titlespacing*{\subsection}{0pt}{0.8em}{0.4em}
\titlespacing*{\subsubsection}{0pt}{0.6em}{0.3em}

% ── Code Listing Style ────────────────────────────────────────────────────────
\lstset{
    backgroundcolor=\color{codebg},
    basicstyle=\ttfamily\small\color{codetext},
    keywordstyle=\color{cyan}\bfseries,
    commentstyle=\color{gray}\itshape,
    stringstyle=\color{yellow},
    breaklines=true,
    frame=none,
    xleftmargin=1em,
    xrightmargin=1em,
    aboveskip=0.5em,
    belowskip=0.5em,
    showstringspaces=false,
    numbers=none,
}

% ── Custom Boxes ──────────────────────────────────────────────────────────────
\tcbset{
    commandbox/.style={
        colback=codebg,
        colframe=primaryblue,
        coltext=codetext,
        fonttitle=\bfseries\color{white},
        attach boxed title to top left={yshift=-2mm, xshift=4mm},
        boxed title style={colback=primaryblue, colframe=primaryblue},
        arc=4pt,
        boxrule=0.8pt,
        left=6pt, right=6pt, top=8pt, bottom=6pt,
        breakable
    },
    infobox/.style={
        colback=lightgray,
        colframe=primaryblue,
        arc=4pt,
        boxrule=0.8pt,
        left=8pt, right=8pt, top=6pt, bottom=6pt
    },
    warningbox/.style={
        colback=mysqlorange!8,
        colframe=mysqlorange,
        arc=4pt,
        boxrule=0.8pt,
        left=8pt, right=8pt, top=6pt, bottom=6pt
    },
    fixbox/.style={
        colback=lightgreen,
        colframe=successgreen,
        arc=4pt,
        boxrule=0.8pt,
        left=8pt, right=8pt, top=6pt, bottom=6pt
    }
}

% ── Header / Footer ───────────────────────────────────────────────────────────
\pagestyle{fancy}
\fancyhf{}
\fancyhead[L]{\small\textcolor{primaryblue}{\textbf{GitDB}} --- Implementation Plan}
\fancyhead[R]{\small\textcolor{gray}{DBMS Mini Project}}
\fancyfoot[C]{\small\textcolor{gray}{\thepage}}
\renewcommand{\headrulewidth}{0.4pt}
\renewcommand{\headrule}{\hbox to\headwidth{\color{primaryblue}\leaders\hrule height \headrulewidth\hfill}}

% ── Document ──────────────────────────────────────────────────────────────────
\begin{document}

% ── Title Page ────────────────────────────────────────────────────────────────
\begin{titlepage}
    \centering
    \vspace*{1.5cm}

    {\Huge\bfseries\color{darkblue} GitDB}\\[0.4em]
    {\Large\color{primaryblue} Implementation Plan}\\[0.2em]
    {\small\color{mysqlorange}\textit{Detailed Phase-by-Phase Build Guide}}\\[0.2em]
    {\small\color{gray}\rule{8cm}{0.5pt}}

    \vspace{1.5cm}

    \begin{tcolorbox}[infobox, width=0.85\textwidth]
        \centering
        {\normalsize\textit{A complete implementation roadmap for GitDB --- covering repo
        scaffold, database setup, snapshot engine, diff engine, two-phase checkout,
        CLI, REST API, and React UI --- with all spec issues identified and fixed.}}
    \end{tcolorbox}

    \vspace{1.5cm}

    {\large\textbf{Institution:} PICT, Pune}\\[0.5em]
    {\large\textbf{Subject:} DISL}\\[0.5em]
    {\large\textbf{Project Type:} Mini Project}\\[0.5em]
    {\large\textbf{Class:} SY-11 \quad|\quad \textbf{Lab Batch:} F-11}\\[0.5em]
    {\large\textbf{Group Size:} 4 Members}\\[1em]

    \vspace{0.5cm}

    \textbf{Team Members:}\\
    \vspace{0.5em}
    \renewcommand{\arraystretch}{1.3}
    \begin{tabular}{>{\centering\arraybackslash}p{1cm} >{\arraybackslash}p{5cm} >{\centering\arraybackslash}p{2.5cm}}
    \toprule
    \textbf{Sr.} & \textbf{Name} & \textbf{Roll No.} \\
    \midrule
    1 & Aditya Malode     & 23325 \\
    2 & Rohan Natekar     & 23330 \\
    3 & Aarush Oak        & 23332 \\
    4 & Prathamesh Kinagi & 23340 \\
    \bottomrule
    \end{tabular}

    \vspace{2.5cm}
    {\small\color{gray} Academic Year 2025--26}
    \vfill
    {\small\color{gray} Compiled on \today}
\end{titlepage}

\tableofcontents
\newpage

% ── Issues Found & Fixed ──────────────────────────────────────────────────────
\section{Issues Found and Fixed in the Spec}

The following issues were identified in the original \texttt{main.tex} and
\texttt{database\_structure.tex} documents. All are corrected in this
implementation plan.

\begin{tcolorbox}[warningbox]
\textbf{Issue 1 --- Password stored in plaintext.}
The spec stores DB credentials only in \texttt{.gitdb/config.json}.
\textbf{Fix:} Use the OS-level \texttt{keyring} library. Only the keyring
service key (e.g.\ \texttt{gitdb:repo\_1}) is stored in config; the actual
password is retrieved at runtime via
\texttt{keyring.get\_password(service, username)}.
\end{tcolorbox}

\begin{tcolorbox}[warningbox]
\textbf{Issue 2 --- No \texttt{UNIQUE} constraint on \texttt{parent\_hash}.}
The spec correctly states no two commits may share a parent (linear chain), but
the DDL never enforces this at the database level.
\textbf{Fix:} Add \texttt{UNIQUE KEY uq\_parent\_hash (parent\_hash)} to the
\texttt{commits} table.
\end{tcolorbox}

\begin{tcolorbox}[warningbox]
\textbf{Issue 3 --- No user registration flow in \texttt{gitdb init}.}
The spec defines a \texttt{user} table with authentication but the \texttt{gitdb
init} command only accepts DB connection arguments --- there is no way to create
a user.
\textbf{Fix:} Add a \texttt{gitdb register} command that creates a user account
with an argon2-hashed password, and make \texttt{--author} required on
\texttt{gitdb init}.
\end{tcolorbox}

\begin{tcolorbox}[warningbox]
\textbf{Issue 4 --- \texttt{gitdb commit --author} is a free-text flag.}
This bypasses the foreign key to \texttt{user} and breaks authorship integrity.
\textbf{Fix:} \texttt{--author} resolves a \texttt{username} from
\texttt{gitdb\_meta.user} via a \texttt{SELECT} query; a raw name string is
never accepted.
\end{tcolorbox}

\begin{tcolorbox}[warningbox]
\textbf{Issue 5 --- \texttt{react-diff-viewer} is unmaintained.}
Last release was 2021 with a peer dependency on React~16 --- incompatible with
the specified React~18.
\textbf{Fix:} Replace with \texttt{@git-diff-view/react}, which is actively
maintained and supports React~18.
\end{tcolorbox}

\begin{tcolorbox}[warningbox]
\textbf{Issue 6 --- No recovery path if Phase~1 checkout fails mid-way.}
If DDL application fails and the subsequent re-apply of the old schema also
fails, the user is left with an inconsistent database and no manual escape route.
\textbf{Fix:} Before any DDL is applied, the pre-checkout snapshot is emitted as
a \texttt{.gitdb/recovery\_\{timestamp\}.sql} file. The user can always run
this file manually to restore the prior state.
\end{tcolorbox}

\begin{tcolorbox}[warningbox]
\textbf{Issue 7 --- \texttt{SELECT *} for snapshotting is non-deterministic.}
The spec describes ``primary-key ordering'' in prose but the SQL shown is a bare
\texttt{SELECT *} with no \texttt{ORDER BY}, making hashes non-reproducible.
\textbf{Fix:} All snapshot queries explicitly use
\texttt{ORDER BY \{pk\_columns\}}.
\end{tcolorbox}

% ── Repo Scaffold ─────────────────────────────────────────────────────────────
\section{Repository Scaffold}

Set up the following directory structure before writing any engine code.
Keeping the CLI, engine, API, and frontend in separate top-level packages
ensures each is independently testable from day one.

\begin{tcolorbox}[commandbox, title=Project Directory Layout]
\begin{lstlisting}
gitdb/
|-- engine/
|   |-- snapshot.py     # DB introspection + row capture
|   |-- diff.py         # Schema + data diff engine
|   |-- hash.py         # SHA-256 Merkle hash
|   `-- checkout.py     # Two-phase checkout
|-- cli/
|   `-- main.py         # click + rich entry point
|-- api/
|   `-- app.py          # Flask REST wrapper
|-- ui/
|   `-- src/            # React 18 + Tailwind + Vite
|-- db/
|   `-- schema.sql      # Canonical DDL (with all fixes applied)
`-- tests/
    |-- unit/
    `-- integration/
\end{lstlisting}
\end{tcolorbox}

\begin{tcolorbox}[infobox]
Use a \texttt{pyproject.toml} with \texttt{pip install -e .} so the
\texttt{gitdb} command is available on PATH during development without manual
path manipulation.
\end{tcolorbox}

% ── Phase 0 ───────────────────────────────────────────────────────────────────
\section{Phase 0 --- Environment Setup}

\begin{itemize}[leftmargin=1.5em, itemsep=0.3em]
    \item \textbf{Runtime requirements:} Python 3.10+, Node 18+, MySQL 8.0+
    \item \textbf{Python dependencies:} \texttt{mysql-connector-python},
    \texttt{sqlglot}, \texttt{click}, \texttt{rich}, \texttt{flask},
    \texttt{flask-cors}, \texttt{argon2-cffi}, \texttt{keyring}
    \item \textbf{Node dependencies (inside \texttt{ui/}):} \texttt{react},
    \texttt{react-dom}, \texttt{@xyflow/react} (react-flow),
    \texttt{@git-diff-view/react}, \texttt{tailwindcss}, \texttt{vite}
    \item Create a local \texttt{gitdb\_meta} MySQL instance for development and
    integration testing
\end{itemize}

% ── Phase 1 ───────────────────────────────────────────────────────────────────
\section{Phase 1 --- \texttt{gitdb\_meta} Database Setup}

Write \texttt{db/schema.sql} with two corrections applied over the original spec.

\begin{tcolorbox}[fixbox]
\textbf{Fix applied:} \texttt{UNIQUE KEY uq\_parent\_hash (parent\_hash)}
added to \texttt{commits} to enforce the linear chain at the database level.
A new \texttt{db\_password\_key VARCHAR(100)} column is added to
\texttt{repository} to store the OS keyring service key --- never the password.
\end{tcolorbox}

\begin{tcolorbox}[commandbox, title=\texttt{db/schema.sql} --- Complete DDL]
\begin{lstlisting}[language=SQL]
CREATE DATABASE IF NOT EXISTS gitdb_meta;
USE gitdb_meta;

-- Users (authentication + authorship)
CREATE TABLE user (
    user_id       INT          NOT NULL AUTO_INCREMENT,
    username      VARCHAR(50)  NOT NULL,
    email         VARCHAR(100) NOT NULL,
    password_hash VARCHAR(255) NOT NULL,
    full_name     VARCHAR(100) NOT NULL,
    is_active     BOOLEAN      NOT NULL DEFAULT TRUE,
    created_at    DATETIME     NOT NULL DEFAULT CURRENT_TIMESTAMP,
    PRIMARY KEY (user_id),
    UNIQUE KEY uq_username (username),
    UNIQUE KEY uq_email    (email)
);

-- Repositories (connection config + HEAD pointer)
CREATE TABLE repository (
    repo_id         INT          NOT NULL AUTO_INCREMENT,
    user_id         INT          NOT NULL,
    repo_name       VARCHAR(100) NOT NULL,
    target_db_name  VARCHAR(100) NOT NULL,
    db_host         VARCHAR(100) NOT NULL,
    db_port         INT          NOT NULL DEFAULT 3306,
    db_user         VARCHAR(50)  NOT NULL,
    db_password_key VARCHAR(100) NULL,   -- keyring key, not the password
    current_hash    CHAR(64)     NULL,
    created_at      DATETIME     NOT NULL DEFAULT CURRENT_TIMESTAMP,
    PRIMARY KEY (repo_id),
    CONSTRAINT fk_repo_user
        FOREIGN KEY (user_id) REFERENCES user (user_id)
);

-- Commits (linear Merkle chain)
CREATE TABLE commits (
    hash         CHAR(64) NOT NULL,
    repo_id      INT      NOT NULL,
    parent_hash  CHAR(64) NULL,
    author_id    INT      NOT NULL,
    message      TEXT     NOT NULL,
    created_at   DATETIME NOT NULL DEFAULT CURRENT_TIMESTAMP,
    PRIMARY KEY (hash),
    UNIQUE KEY uq_parent_hash (parent_hash),  -- enforces linear chain
    CONSTRAINT fk_commit_repo
        FOREIGN KEY (repo_id)     REFERENCES repository (repo_id),
    CONSTRAINT fk_commit_parent
        FOREIGN KEY (parent_hash) REFERENCES commits    (hash),
    CONSTRAINT fk_commit_author
        FOREIGN KEY (author_id)   REFERENCES user       (user_id)
);

-- Deferred FK: resolves circular dependency between repository and commits
ALTER TABLE repository
    ADD CONSTRAINT fk_repo_head
        FOREIGN KEY (current_hash) REFERENCES commits (hash);

-- Snapshots (full schema + data per table per commit)
CREATE TABLE snapshots (
    id           BIGINT       NOT NULL AUTO_INCREMENT,
    commit_hash  CHAR(64)     NOT NULL,
    table_name   VARCHAR(255) NOT NULL,
    ddl_json     LONGTEXT     NOT NULL,
    rows_json    LONGTEXT     NOT NULL,
    row_count    INT          NOT NULL,
    captured_at  DATETIME     NOT NULL DEFAULT CURRENT_TIMESTAMP,
    PRIMARY KEY (id),
    UNIQUE KEY uq_commit_table (commit_hash, table_name),
    CONSTRAINT fk_snap_commit
        FOREIGN KEY (commit_hash) REFERENCES commits (hash)
);
\end{lstlisting}
\end{tcolorbox}

% ── Phase 2 ───────────────────────────────────────────────────────────────────
\section{Phase 2 --- \texttt{gitdb register} and \texttt{gitdb init}}

\subsection{\texttt{gitdb register} (new command)}

This command is required because \texttt{gitdb init} needs a valid
\texttt{author\_id} FK, but the original spec provides no mechanism to create
a user.

\begin{tcolorbox}[commandbox, title=\texttt{gitdb register}]
\begin{lstlisting}
gitdb register
\end{lstlisting}
\textcolor{codetext}{\small Prompts interactively for \texttt{username},
\texttt{email}, \texttt{full\_name}, and \texttt{password} (via
\texttt{getpass} --- never echoed). Hashes the password with
\texttt{argon2-cffi} (\texttt{argon2.PasswordHasher().hash(password)}).
Inserts into \texttt{gitdb\_meta.user} and prints the new \texttt{user\_id}.}
\end{tcolorbox}

\subsection{\texttt{gitdb init}}

\begin{tcolorbox}[commandbox, title=\texttt{gitdb init}]
\begin{lstlisting}
gitdb init --host localhost --port 3306      \
           --user root --password secret     \
           --database myapp_db              \
           --repo-name my-repo              \
           --author <username>
\end{lstlisting}
\textcolor{codetext}{\small Connects to MySQL, creates \texttt{gitdb\_meta}
and all four tables if not present (runs \texttt{schema.sql}). Validates that
\texttt{--author} resolves to an active user in \texttt{gitdb\_meta.user}.
Inserts a row into \texttt{repository} with \texttt{current\_hash = NULL}.
Stores the password securely via
\texttt{keyring.set\_password("gitdb", "repo\_\{id\}", password)} and writes
only the keyring service key to \texttt{repository.db\_password\_key}.
Writes \texttt{.gitdb/config.json} with no password.}
\end{tcolorbox}

\begin{tcolorbox}[commandbox, title=\texttt{.gitdb/config.json} --- written by \texttt{gitdb init}]
\begin{lstlisting}[language={}]
{
  "repo_id": 1,
  "host":    "localhost",
  "port":    3306,
  "db_user": "root",
  "db_name": "myapp_db"
}
\end{lstlisting}
\textcolor{codetext}{\small No password is ever written to this file.}
\end{tcolorbox}

% ── Phase 3 ───────────────────────────────────────────────────────────────────
\section{Phase 3 --- Snapshot Engine (\texttt{engine/snapshot.py})}

Core data capture layer. Called at every \texttt{gitdb commit}.

\begin{tcolorbox}[commandbox, title=Function Signature]
\begin{lstlisting}[language=Python]
def capture_snapshot(conn, db_name: str) -> dict[str, TableSnapshot]
\end{lstlisting}
\end{tcolorbox}

\subsection{Capture Steps}

\begin{enumerate}[leftmargin=1.5em, itemsep=0.3em]
    \item Query \texttt{information\_schema.TABLES} to get all
    \texttt{BASE TABLE} names for \texttt{db\_name}.

    \item For each table, check primary key existence via
    \texttt{information\_schema.KEY\_COLUMN\_USAGE WHERE CONSTRAINT\_NAME =
    'PRIMARY'}. If none exists: emit a \texttt{rich} warning panel and skip
    data capture for that table (\texttt{rows\_json = []},
    \texttt{row\_count = 0}). DDL is still captured.

    \item \textbf{Capture \texttt{ddl\_json}} --- structured representation
    from \texttt{information\_schema.COLUMNS} plus the raw output of
    \texttt{SHOW CREATE TABLE} stored under the key \texttt{raw\_ddl}.
    The \texttt{raw\_ddl} string is used by \texttt{sqlglot} in the diff engine.

    \item \textbf{Capture \texttt{rows\_json}} --- \texttt{ORDER BY} on all
    primary key columns is mandatory for deterministic hashing:
\end{enumerate}

\begin{tcolorbox}[commandbox, title=Row Capture Query (deterministic ordering)]
\begin{lstlisting}[language=SQL]
SELECT * FROM `{table}` ORDER BY `pk_col1`, `pk_col2`;
\end{lstlisting}
\textcolor{codetext}{\small If \texttt{COUNT(*) > 10000}, raise
\texttt{SnapshotRowLimitError} with a clear message before fetching any rows.}
\end{tcolorbox}

\subsection{\texttt{TableSnapshot} Dataclass --- Shared Integration Contract}

\begin{tcolorbox}[fixbox]
\textbf{Integration contract:} Aditya defines and exports this dataclass from
\texttt{engine/snapshot.py} in \textbf{Week~1}. All other team members import
it. This is the single shared type that prevents workstreams from blocking
each other.
\end{tcolorbox}

\begin{tcolorbox}[commandbox, title=\texttt{TableSnapshot} Dataclass]
\begin{lstlisting}[language=Python]
@dataclass
class TableSnapshot:
    table_name: str
    ddl_json:   dict        # column metadata + raw_ddl string
    rows_json:  list[dict]  # list of row dicts, ordered by PK
    row_count:  int
    pk_columns: list[str]   # empty list if no primary key
\end{lstlisting}
\end{tcolorbox}

% ── Phase 4 ───────────────────────────────────────────────────────────────────
\section{Phase 4 --- Hashing Engine (\texttt{engine/hash.py})}

\begin{tcolorbox}[commandbox, title=Function Signature]
\begin{lstlisting}[language=Python]
def compute_commit_hash(
    snapshot: dict[str, TableSnapshot],
    parent_hash: str | None
) -> str
\end{lstlisting}
\end{tcolorbox}

\begin{enumerate}[leftmargin=1.5em, itemsep=0.3em]
    \item Serialise \texttt{snapshot} to JSON with sorted keys:
    \texttt{json.dumps(snapshot\_dict, sort\_keys=True)}
    \item Concatenate with parent hash (empty string for the root commit)
    \item Return \texttt{hashlib.sha256(payload).hexdigest()} --- 64-char hex
\end{enumerate}

\begin{tcolorbox}[commandbox, title=Hash Construction]
\begin{lstlisting}[language=Python]
payload = json_bytes + (parent_hash or "").encode("utf-8")
return hashlib.sha256(payload).hexdigest()
\end{lstlisting}
\end{tcolorbox}

\begin{tcolorbox}[infobox]
\textbf{Determinism guarantee:} Same DB state + same parent hash always
produces the same commit hash. Any modification to a historical snapshot
immediately invalidates all subsequent hashes in the Merkle chain.
\end{tcolorbox}

% ── Phase 5 ───────────────────────────────────────────────────────────────────
\section{Phase 5 --- Diff Engine (\texttt{engine/diff.py})}

Two independent sub-engines, composed into a single \texttt{diff(hash1, hash2)}
function returning a \texttt{DiffResult}.

\subsection{Schema Diff (using \texttt{sqlglot})}

\begin{enumerate}[leftmargin=1.5em, itemsep=0.3em]
    \item Load \texttt{ddl\_json["raw\_ddl"]} for both snapshots from
    \texttt{gitdb\_meta.snapshots}
    \item Parse each DDL string:
    \texttt{sqlglot.parse\_one(raw\_ddl, dialect="mysql")}
    \item Compare column sets and generate SQL accordingly
\end{enumerate}

\vspace{0.5em}
\renewcommand{\arraystretch}{1.4}
\begin{tabularx}{\textwidth}{>{\ttfamily}p{4.5cm} X}
\rowcolor{tableheader}
\textcolor{white}{\normalfont\bfseries Scenario} &
\textcolor{white}{\normalfont\bfseries Generated Output} \\
\midrule
Table in new, not in old      & \texttt{CREATE TABLE ...} \\
\rowcolor{lightgray}
Table in old, not in new      & \texttt{DROP TABLE IF EXISTS ...} \\
Column added                   & \texttt{ALTER TABLE ... ADD COLUMN ...} \\
\rowcolor{lightgray}
Column dropped                 & \texttt{ALTER TABLE ... DROP COLUMN ...} \\
Column type changed            & Warning banner + SQL comment (not silently handled) \\
\rowcolor{lightgray}
Table or column renamed        & Warning banner + SQL comment (ambiguous --- user must resolve) \\
\end{tabularx}
\vspace{0.5em}

\begin{tcolorbox}[infobox]
Column type changes and renames are intentionally flagged rather than
silently auto-handled, consistent with the design philosophy of
Prisma and Atlas: correctness over convenience.
\end{tcolorbox}

\subsection{Data Diff (Primary-Key Anchored, $O(n)$)}

\begin{enumerate}[leftmargin=1.5em, itemsep=0.3em]
    \item Load \texttt{rows\_json} for both snapshots from
    \texttt{gitdb\_meta.snapshots}
    \item Index both as \texttt{dict[pk\_tuple, row\_dict]}
    \item Compare keys and values:
\end{enumerate}

\vspace{0.5em}
\renewcommand{\arraystretch}{1.4}
\begin{tabularx}{\textwidth}{>{\ttfamily}p{5cm} X}
\rowcolor{tableheader}
\textcolor{white}{\normalfont\bfseries Scenario} &
\textcolor{white}{\normalfont\bfseries Generated Output} \\
\midrule
Key in new, not in old         & \texttt{INSERT INTO ... VALUES (...)} \\
\rowcolor{lightgray}
Key in both, values differ     & \texttt{UPDATE ... SET ... WHERE pk = ...} \\
Key in old, not in new         & \texttt{DELETE FROM ... WHERE pk = ...} \\
\end{tabularx}
\vspace{0.5em}

\subsection{\texttt{DiffResult} Output Structure}

\begin{tcolorbox}[commandbox, title=\texttt{DiffResult} Dataclass]
\begin{lstlisting}[language=Python]
@dataclass
class DiffResult:
    schema_sql: list[str]   # ordered: CREATE, ALTER, DROP
    data_sql:   list[str]   # ordered: INSERT, UPDATE, DELETE
    warnings:   list[str]   # type changes, renames, no-PK tables
\end{lstlisting}
\end{tcolorbox}

% ── Phase 6 ───────────────────────────────────────────────────────────────────
\section{Phase 6 --- Checkout Engine (\texttt{engine/checkout.py})}

Two-phase strategy that handles MySQL DDL auto-commit explicitly.

\begin{tcolorbox}[warningbox]
\textbf{MySQL DDL Note:} \texttt{ALTER TABLE}, \texttt{DROP TABLE}, and
\texttt{CREATE TABLE} issue implicit commits and cannot be rolled back in a
standard transaction. GitDB's two-phase checkout accounts for this explicitly.
\end{tcolorbox}

\subsection{Phase 1 --- Schema}

\begin{tcolorbox}[commandbox, title=Phase 1 --- Schema Application with Recovery]
\begin{lstlisting}[language=Python]
for stmt in diff.schema_sql:
    try:
        cursor.execute(stmt)
    except Exception as e:
        # DDL already auto-committed -- restore from snapshot
        restore_schema_from_snapshot(conn, old_snapshot)
        write_recovery_file(
            old_snapshot,
            path=".gitdb/recovery_{timestamp}.sql"
        )
        raise CheckoutSchemaError(str(e))
\end{lstlisting}
\textcolor{codetext}{\small \texttt{restore\_schema\_from\_snapshot}
re-executes the \texttt{raw\_ddl} strings from the pre-checkout snapshot.
\texttt{write\_recovery\_file} emits a \texttt{.sql} file the user can
run manually if restore also fails --- this is the fix for Issue~6.}
\end{tcolorbox}

\subsection{Phase 2 --- Data}

\begin{tcolorbox}[commandbox, title=Phase 2 --- Data Patch (Transactional)]
\begin{lstlisting}[language=SQL]
SET FOREIGN_KEY_CHECKS = 0;
-- ... execute all INSERT / UPDATE / DELETE statements ...
SET FOREIGN_KEY_CHECKS = 1;
\end{lstlisting}
\textcolor{codetext}{\small Wrapped in an explicit \texttt{BEGIN} /
\texttt{COMMIT} transaction. On any exception, \texttt{ROLLBACK} is called
automatically. Since the target state is always a known-good snapshot,
re-enabling foreign key checks after the patch guarantees referential
integrity of the final state.}
\end{tcolorbox}

\subsection{On Success}

\begin{tcolorbox}[commandbox, title=HEAD Pointer Update]
\begin{lstlisting}[language=SQL]
UPDATE gitdb_meta.repository
SET    current_hash = %s
WHERE  repo_id = %s;
\end{lstlisting}
\end{tcolorbox}

% ── Phase 7 ───────────────────────────────────────────────────────────────────
\section{Phase 7 --- CLI (\texttt{cli/main.py})}

Built with \texttt{click} and \texttt{rich}. All commands load
\texttt{.gitdb/config.json}, retrieve the password from the OS keyring,
and open two MySQL connections: one to \texttt{gitdb\_meta}, one to the
target database.

\subsection{Command Reference}

\begin{tcolorbox}[commandbox, title=\texttt{gitdb register}]
\begin{lstlisting}
gitdb register
\end{lstlisting}
\textcolor{codetext}{\small Creates a user in \texttt{gitdb\_meta.user}
with an argon2-hashed password. Prompts interactively --- no plaintext
password is ever passed as a CLI argument.}
\end{tcolorbox}

\begin{tcolorbox}[commandbox, title=\texttt{gitdb init}]
\begin{lstlisting}
gitdb init --host HOST --port PORT --user USER \
           --password PASS --database DB       \
           --repo-name NAME --author USERNAME
\end{lstlisting}
\textcolor{codetext}{\small Creates \texttt{gitdb\_meta}, registers the
repository, writes \texttt{.gitdb/config.json}, and stores the password
in the OS keyring.}
\end{tcolorbox}

\begin{tcolorbox}[commandbox, title=\texttt{gitdb commit}]
\begin{lstlisting}
gitdb commit -m "message" --author USERNAME
\end{lstlisting}
\textcolor{codetext}{\small Resolves \texttt{--author} to a \texttt{user\_id}
via \texttt{SELECT user\_id FROM gitdb\_meta.user WHERE username = \%s AND
is\_active = TRUE}. Aborts with a clear error if not found. Calls
\texttt{capture\_snapshot} $\to$ \texttt{compute\_commit\_hash} $\to$ inserts
into \texttt{commits} and \texttt{snapshots} $\to$ updates
\texttt{current\_hash}.}
\end{tcolorbox}

\begin{tcolorbox}[commandbox, title=\texttt{gitdb log}]
\begin{lstlisting}
gitdb log [--oneline] [--graph]
\end{lstlisting}
\textcolor{codetext}{\small Walks the \texttt{parent\_hash} chain from
\texttt{current\_hash}. Renders with \texttt{rich.table}.
\texttt{--oneline} shows \texttt{<short\_hash> <message>}.
\texttt{--graph} renders a simple ASCII \texttt{*} chain.}
\end{tcolorbox}

\begin{tcolorbox}[commandbox, title=\texttt{gitdb diff}]
\begin{lstlisting}
gitdb diff <hash1> <hash2> [--schema-only] [--data-only]
\end{lstlisting}
\textcolor{codetext}{\small Runs the diff engine. Prints generated SQL
with \texttt{rich.syntax} (SQL highlighting). Warnings shown as amber
panels above the output.}
\end{tcolorbox}

\begin{tcolorbox}[commandbox, title=\texttt{gitdb checkout}]
\begin{lstlisting}
gitdb checkout <commit_hash>
\end{lstlisting}
\textcolor{codetext}{\small Executes the two-phase checkout. On Phase~1
failure: prints the recovery file path. On Phase~2 failure: shows rollback
confirmation.}
\end{tcolorbox}

\begin{tcolorbox}[commandbox, title=\texttt{gitdb status}]
\begin{lstlisting}
gitdb status
\end{lstlisting}
\textcolor{codetext}{\small Compares the live database against the HEAD
snapshot. Reports: added tables (in live DB, not in snapshot), dropped tables
(in snapshot, not in live DB), modified tables (column differences), and row
count deltas per table.}
\end{tcolorbox}

% ── Phase 8 ───────────────────────────────────────────────────────────────────
\section{Phase 8 --- Flask REST API (\texttt{api/app.py})}

Thin HTTP wrapper. All business logic lives in the engine layer --- Flask only
calls it and serialises results to JSON. Enable \texttt{flask-cors} for local
development (React runs on a different port).

\vspace{0.5em}
\renewcommand{\arraystretch}{1.4}
\begin{tabularx}{\textwidth}{>{\ttfamily}p{4.5cm} >{\ttfamily}p{1.5cm} X}
\rowcolor{tableheader}
\textcolor{white}{\normalfont\bfseries Endpoint} &
\textcolor{white}{\normalfont\bfseries Method} &
\textcolor{white}{\normalfont\bfseries Description} \\
\midrule
/commits & GET & JSON array of commits, reverse-chronological \\
\rowcolor{lightgray}
/commit  & POST & Body: \{ message, author\_username \}. Returns \{ hash \} \\
/diff/<h1>/<h2> & GET & Returns \{ schema\_sql, data\_sql, warnings \} \\
\rowcolor{lightgray}
/checkout/<hash> & POST & Two-phase checkout. Returns \{ ok \} or \{ error, recovery\_file \} \\
/status & GET & Returns \{ added\_tables, dropped\_tables, modified\_tables, row\_deltas \} \\
\end{tabularx}
\vspace{0.5em}

\begin{tcolorbox}[infobox]
All error responses return \texttt{\{ "error": "..." \}} with an appropriate
HTTP status code. The \texttt{recovery\_file} field is included in checkout
error responses so the UI can surface it to the user directly.
\end{tcolorbox}

% ── Phase 9 ───────────────────────────────────────────────────────────────────
\section{Phase 9 --- React UI (\texttt{ui/})}

Vite + React~18 + Tailwind CSS. Three views navigated from a persistent
sidebar.

\subsection{View 1 --- Commit Graph}

Built with \texttt{react-flow}. Each commit is rendered as a node in a
vertical linear chain.

\begin{itemize}[leftmargin=1.5em, itemsep=0.3em]
    \item Node content: short hash (first 8 characters), commit message,
    author \texttt{full\_name}, timestamp
    \item HEAD node visually highlighted with an accent border and background
    \item Clicking a node opens a side panel showing the full hash, a checkout
    button, and a ``diff from HEAD'' shortcut
\end{itemize}

\subsection{View 2 --- Diff Viewer}

\begin{itemize}[leftmargin=1.5em, itemsep=0.3em]
    \item Two commit dropdowns to select \texttt{hash1} and \texttt{hash2}
    \item Calls \texttt{GET /diff/<h1>/<h2>}
    \item Renders using \texttt{@git-diff-view/react} (replaces the
    unmaintained \texttt{react-diff-viewer})
    \item Schema SQL in left panel, data SQL in right panel
    \item Toggle: Schema only / Data only / Both
    \item Warnings (column type changes, no-PK tables) shown as amber banners
    above the diff output
\end{itemize}

\subsection{View 3 --- Schema Browser}

\begin{itemize}[leftmargin=1.5em, itemsep=0.3em]
    \item Commit selector dropdown
    \item Displays \texttt{ddl\_json} for each table at the selected commit
    as a structured table with columns: name, type, nullable, key, default,
    extra
    \item Read-only --- no mutations from this view
\end{itemize}

% ── Phase 10 ──────────────────────────────────────────────────────────────────
\section{Phase 10 --- Testing Strategy}

\subsection{Unit Tests (pytest, no live DB required)}

\begin{itemize}[leftmargin=1.5em, itemsep=0.3em]
    \item \texttt{test\_diff.py} --- feed two hand-crafted
    \texttt{TableSnapshot} dicts, assert exact SQL output for adds, drops,
    modifies, and deletes
    \item \texttt{test\_hash.py} --- assert determinism (same input $\to$
    same hash; different input $\to$ different hash; parent hash propagation)
    \item \texttt{test\_snapshot.py} --- mock \texttt{mysql-connector-python}
    cursor, assert correct queries are issued including \texttt{ORDER BY pk}
\end{itemize}

\subsection{Integration Tests (require local MySQL instance)}

\begin{itemize}[leftmargin=1.5em, itemsep=0.3em]
    \item \textbf{Full round-trip:} commit state~A $\to$ modify DB $\to$
    commit state~B $\to$ \texttt{checkout A} $\to$ assert DB matches
    state~A exactly
    \item \textbf{FK ordering:} snapshot a DB with FK-linked tables,
    checkout, assert no referential integrity violations
    \item \textbf{Row limit:} seed 10,001 rows, assert
    \texttt{SnapshotRowLimitError} is raised before any data is fetched
    \item \textbf{Schema-only checkout:} add a column $\to$ commit $\to$
    drop it $\to$ commit $\to$ checkout prior commit $\to$ assert column
    restored correctly
\end{itemize}

\subsection{API Tests (pytest + \texttt{flask.testing.FlaskClient})}

\begin{itemize}[leftmargin=1.5em, itemsep=0.3em]
    \item Happy path + error path for each endpoint
    \item Checkout failure path: assert \texttt{recovery\_file} field is
    present in the error response JSON
\end{itemize}

% ── Phase 11 ──────────────────────────────────────────────────────────────────
\section{Phase 11 --- Team Ownership and Delivery Schedule}

\renewcommand{\arraystretch}{1.5}
\begin{tabularx}{\textwidth}{>{\bfseries}p{3.8cm} X >{\centering\arraybackslash}p{2cm}}
\rowcolor{tableheader}
\textcolor{white}{\normalfont\bfseries Member} &
\textcolor{white}{\normalfont\bfseries Owns} &
\textcolor{white}{\normalfont\bfseries By} \\
\midrule
Aditya Malode
    & \texttt{engine/snapshot.py}, \texttt{engine/hash.py},
      \texttt{db/schema.sql}, \texttt{gitdb init},
      \texttt{gitdb register}
    & Week 1 \\
\rowcolor{lightgray}
Rohan Natekar
    & \texttt{engine/diff.py} (schema diffing with \texttt{sqlglot}),
      CLI: \texttt{gitdb diff}, \texttt{gitdb status}
    & Week 2 \\
Aarush Oak
    & \texttt{engine/checkout.py} (two-phase),
      CLI: \texttt{gitdb commit}, \texttt{gitdb log},
      \texttt{gitdb checkout}
    & Week 2 \\
\rowcolor{lightgray}
Prathamesh Kinagi
    & \texttt{api/app.py} (Flask REST API),
      entire React UI (all three views)
    & Week 3 \\
\end{tabularx}

\vspace{1em}

\begin{tcolorbox}[fixbox]
\textbf{Integration contract:} Aditya defines and exports the
\texttt{TableSnapshot} dataclass from \texttt{engine/snapshot.py} in
\textbf{Week~1}. All other team members import it from this module. This is
the single shared type that keeps all four workstreams from blocking each
other.
\end{tcolorbox}

\begin{tcolorbox}[commandbox, title=\texttt{TableSnapshot} --- Week 1 Contract (engine/snapshot.py)]
\begin{lstlisting}[language=Python]
@dataclass
class TableSnapshot:
    table_name: str
    ddl_json:   dict        # column metadata + raw_ddl string
    rows_json:  list[dict]  # list of row dicts, ordered by PK
    row_count:  int
    pk_columns: list[str]   # empty list if no primary key
\end{lstlisting}
\end{tcolorbox}

\subsection{Dependency Map}

\begin{tcolorbox}[infobox]
\begin{itemize}[leftmargin=1.5em, itemsep=0.2em]
    \item \texttt{gitdb init / register} $\to$ \texttt{schema.sql}
    (Phase~1) + \texttt{keyring} (Phase~2)
    \item \texttt{gitdb commit} $\to$ \texttt{snapshot.py} (Phase~3)
    + \texttt{hash.py} (Phase~4)
    \item \texttt{gitdb diff} $\to$ \texttt{diff.py} (Phase~5)
    \item \texttt{gitdb checkout} $\to$ \texttt{diff.py} (Phase~5)
    + \texttt{checkout.py} (Phase~6)
    \item Flask API $\to$ all engine modules
    \item React UI $\to$ Flask API
\end{itemize}
\end{tcolorbox}

\end{document}